\icode{N22E-A-SI-L2-260930-008}
\nceeID[4]{2022湖北}{1}
上世纪四十年代初，我国科学家王淦昌先生首先提出证明中微子存在的实验方案：如果静止原子核\ce{^7_4Be}俘获核外\(K\)层电子\(e\)，可生成一个新原子核\(X\)，并放出中微子\(v_e\)，即\ce{^7_4Be + ^0_{-1}e \to ^7_3X + ^0_0v_e}。根据核反应后原子核\(X\)的动能和动量，可以间接测量中微子的能量和动量，进而确定中微子的存在。下列说法正确的是\paren[]
\begin{choices}
	\item 原子核\(X\)是\ce{^7_3Li}
	\item 核反应前后的总质子数不变
	\item 核反应前后总质量数不同
	\item 中微子\(v_e\)的电荷量与电子的相同
\end{choices}
\begin{answer}
	A
\end{answer}